\section{مقارنة استراتيجيات التحكم واختيار الخوارزمية الأمثل}
\label{sec:control_comparison}

بعد التطبيق والتحليل المعمق لخوارزميات التحكم الثلاث (\LR{PID}، \LR{SMC}، و \LR{ACTC}) على النموذج الديناميكي الشامل لروبوت (\LR{SCARA})، واختبارها تحت ظروف تحميل متغيرة (من حالة اللاحمل إلى حمل أقصى $m_L = 10 \text{ kg}$)، يمكن تلخيص الفروقات الجوهرية بين هذه الاستراتيجيات من حيث الأداء الديناميكي، المتانة، والعبء الحسابي في الجدول التالي:
\begin{table}[H]
    \centering
    \caption{مقارنة هندسية بين استراتيجيات التحكم المطبقة على الروبوت}
    \label{tab:control_comparison}
    % إضافة تعليمات إيقاف المط (raggedright) والتوسيط (centering) للأعمدة
    \begin{tabular}{|>{\centering\arraybackslash}p{0.14\linewidth}|>{\raggedright\arraybackslash}p{0.21\linewidth}|>{\raggedright\arraybackslash}p{0.22\linewidth}|>{\raggedright\arraybackslash}p{0.21\linewidth}|>{\centering\arraybackslash}p{0.12\linewidth}|}
        \hline
        \textbf{الخوارزمية} & \textbf{دقة التتبع الموضعي} & \textbf{المتانة ضد تغيرات الحمل ($m_L$)} & \textbf{زمن الاستقرار} & \textbf{التعقيد الحسابي برمجياً} \\
        \hline
        \textbf{الكلاسيكي (\LR{PID})} & ممتازة في حالة اللاحمل، وضعيفة عند السرعات العالية. & منهارة (تأثر شديد بزيادة عزم القصور الذاتي للوصلة). & سريع في اللاحمل، وينهار (يفقد الاستقرار) عند إضافة الحمل. & بسيط جداً (لا يتطلب نموذج رياضي). \\
        \hline
        \textbf{الانزلاقي (\LR{SMC})} & عالية جداً. & ممتازة (مناعة تامة ضد الاضطرابات الخارجية). & \centering{جيد} & متوسط (يتطلب معرفة النموذج الاسمي). \\
        \hline
        \textbf{التكيفي (\LR{ACTC})} & عالية جداً. & ممتازة (تكيف ذاتي وتقدير لحظي ). & \centering{قياسي وسريع } & عالي (يتطلب تفكيك وحساب مصفوفة الانحدار $Y$). \\
        \hline
    \end{tabular}
\end{table}

\subsection*{القرار الهندسي: اعتماد خوارزمية (\LR{ACTC})}
بناءً على المعطيات التحليلية السابقة وطبيعة المهمة الصناعية لروبوت (\LR{SCARA}) المدروس (والتي تتمثل في نقل حمولات متغيرة الوزن بسرعات عالية)، \textbf{تم اختيار خوارزمية التحكم التكيفي المعتمد على العزم المحسوب (\LR{ACTC})} لتكون الاستراتيجية النهائية والمعتمدة في هذا المشروع. 

يأتي هذا الاختيار مبرراً بالنقاط والمزايا الهندسية التالية:
\begin{enumerate}
    \item \textbf{زمن استجابة قياسي:} أثبتت الخوارزمية قدرتها على استعادة استقرار النظام وتصفير خطأ التتبع خلال زمن قصير جداً ($0.5 \text{ s}$) حتى عند تعرض الروبوت  لصدمة تحميل ($10 \text{ kg}$)، مع امتصاص ذكي لعزوم التداخل الديناميكي بين المفاصل.
    \item \textbf{الاستغناء عن المعايرة اليدوية:} قدرة المتحكم التلقائية على تقدير الزيادة في القصور الذاتي وتحديث المعاملات المجهولة (مثل $\hat{\theta}_{M,3}$) لحظياً، مما يغني عن الحاجة لإعادة ضبط أرباح التحكم يدوياً عند كل تغيير في وزن القطعة المنقولة.
    \item \textbf{الحفاظ على العمر الافتراضي للمحركات:} على عكس متحكم (\LR{SMC}) الذي يضمن المتانة على حساب إرهاق المحركات بجهود تبديلية عالية التردد (\LR{Chattering})، توفر خوارزمية (\LR{ACTC}) إشارات تحكم ناعمة ومستمرة خالية من الاهتزازات، مما يقلل من التآكل الميكانيكي في المفاصل والارتفاع الحراري في ملفات المحركات.
\end{enumerate}

بهذا الاستنتاج، نكون قد أتممنا تصميم واختبار "العقل" البرمجي القادر على قيادة ذراع الروبوت بكفاءة صناعية عالية.